% Joint continuation increment, v0.4, 2026-11-07.
% Standalone module: requires amsmath, amssymb, amsthm and the usual theorem
% environments. It does not redefine the report-row error parameter.
% Complete proof record: joint_continuation.md, same directory.
\subsection{Joint continuation from an epoch root to a message}
\label{sec:joint-continuation}

A recovered root feeds confirmation, traffic derivation and the next epoch.
We treat these uses jointly and then isolate a selected message key.
The resulting continuation permits later contributions, successor roots and
the next chain seed to be disclosed. The current root is an internal value:
all outputs are regenerated in each hybrid, rather than fixed before a
root replacement.

\paragraph{Primitive experiments.}
Let $F(k,m)$ be HMAC-SHA256 over one public quantum compression oracle
$h:\{0,1\}^{512}\times\{0,1\}^{256}\to\{0,1\}^{256}$, in message-block,
chaining-state order. All keys below have 32 bytes. Put
\[
 \Gamma(a)=\bigl(h(I(a),\mathrm{IV}),h(O(a),\mathrm{IV})\bigr),
\]
where $I(a)=(a\Vert0^{32})\oplus\mathtt{36}^{64}$ and
$O(a)=(a\Vert0^{32})\oplus\mathtt{5c}^{64}$.
Initial auxiliary state is independent of $h$ and the designated hidden
uniform key. All oracle-dependent preparation is executed and counted,
together with the complete continuation.

Write $\varepsilon_{\rm Hjoint}(C,T,q,L)$ for an output-bit gap bound between
\[
 \operatorname{Col}(a)=F(a,R),\quad
 \operatorname{Row}(m)=F(R,m)
 \quad\hbox{and}\quad
 V(\Gamma(a)),\ U(m),
\]
where $R$ is hidden uniform, $U,V$ are independent private random functions,
there are at most $C$ classical column calls, $T$ classical row calls of
length at most $L$, and $q$ ordinary compression queries, which may be
quantum. Equal inputs are cached. Distinct salts with equal $\Gamma$ share
their column reply. The ideal here is the independent-table ideal, without
a separately imposed cross-pair repair.

Write $\varepsilon_{\rm Hrow}(T,q,L)$ for the corresponding real-HMAC-row
versus private-random-function bound for one uniform hidden key and
\emph{multiple adaptive messages}. This differs from the immutable report
row experiment. Finally, $\varepsilon_{\rm Hcol}(q)$ bounds one distinct
zero-salt evaluation $F(0,S)$, for hidden uniform $S$, versus an independent
uniform reply; exact repeats reuse it. The retained compression-QROM
arguments give admissible bounds for these operational experiments.

\paragraph{The finite continuation.}
The initial root $R$ is sampled independently of $h$ and the initial
auxiliary state, before any use. It may be presampled and withheld during
counted preparation, which chooses the public session identifier $s$ and
transcript inputs. The previous root and old wrapping keys may be known.
The current root has only the uses
\begin{equation}
 P_0=F(0,R),\qquad P_s=F(s,R),\qquad P_x^+=F(R,S_x^+).
 \label{eq:joint-root-uses}
\end{equation}
There are at most $t_+$ distinct 32-byte successor contributions $S_x^+$.
They are supplied by known local computation, may be disclosed, and may
depend on the available view. The controller may compute and disclose
\[
 R_x^+=\operatorname{Expand}(P_x^+,I_x^+,32)
\]
and run arbitrary known subsequent computation from these values with the
same public oracle. Repeated arguments remain consistent.

Let $D$ and $D_M$ denote the actual 40-byte rekey and 43-byte message domains.
The three confirmation information strings are
\[
\begin{aligned}
 I_r&=D\Vert\texttt{responder-confirmation},\\
 I_f&=D\Vert\texttt{final-confirmation},\\
 I_a&=D\Vert\texttt{receipt-confirmation}.
\end{aligned}
\]
Their keys are $\operatorname{Expand}(P_0,I_d,32)$. The controller uses
these keys in the real MAC, verification and signing algorithms.
At most $M\ge1$ distinct traffic information strings have the native form
\begin{equation}
\begin{split}
 I_T={}&D\Vert\texttt{epoch-traffic/HKDF-SHA256/ChaCha20Poly1305/}\\
 &\Vert\mathit{context}_{32}\Vert\operatorname{BE64}(e)
 \Vert\mathit{transcriptDigest}_{32}.
\end{split}
 \label{eq:joint-traffic-info}
\end{equation}
It may be chosen after the response confirmation and its signature: the
native transcript digest hashes the complete offer and response wires.
Every repeated complete $I_T$ reuses its entry, regardless of endpoint or
wire representation. At the first allocation of one new string $I^*$,
mark that entry and a direction $d^*\in\{1,2\}$ before any of its blocks
are produced. The choice may use the available prefix; a previously
allocated entry cannot be reclassified as fresh. Later equal-info
aliases inherit the same protection. This is a profile for a designated
fresh traffic allocation, without a cross-endpoint wire-uniqueness premise.

The four 32-byte output blocks are
\[
 (s_{I,1},s_{I,2},a_{I,1},a_{I,2})
 =\operatorname{Expand}(P_s,I,128).
\]
The selected seed $s^*=s_{I^*,d^*}$, designated before allocation,
has only the selected first-step use
\begin{equation}
 \begin{split}
 I_M&=D_M\Vert\texttt{chain}\Vert\operatorname{BE64}(0),\\
 P_M&=F(0,s^*),\qquad
 (s',K)=\operatorname{Expand}(P_M,I_M,64).
 \end{split}
 \label{eq:joint-chain-step}
\end{equation}
The opposite seed, confirmation keys, ACK keys, blocks at every distinct
nonmarked information string and next seed $s'$ may be disclosed.
A repeated $I^*$ remains the same protected entry.
Real ACK computation and arbitrary candidate-ACK verification
remain available. Subsequent chain computations use $s'$ honestly.
The selected key $K$ has only its declared application encryption/decryption
use. Current-root, PRK, old selected-seed and message-key bytes have no
separate public return, comparison or hashing interface. Private persistence
uses an already justified opaque interface; original candidate-local stops
and projection accounting remain in force.

\begin{lemma}[Adaptive information and fresh traffic blocks]
\label{lem:joint-expand-blocks}
Replace each PRK-class HMAC row by an independent private random function,
available only through HKDF-Expand. The resulting interface is exactly an
independent block table indexed by the class, complete information string
and block index, with consistent prefixes. This holds for adaptively
chosen information strings. In particular, the first four traffic blocks
are fresh at installation despite the confirmation dependence of $I_T$.
\end{lemma}
\begin{proof}
For one private row $W$, the inputs are
$I\Vert1$ and $T_{I,j-1}\Vert I\Vert j$ for $j\ge2$, where each preceding
block has exactly 32 bytes and the counter has one byte. Different counters
separate inputs. For the same counter, deleting the counter and, when
present, the fixed-width preceding block shows that equality requires equal
$I$. This is true for every preceding output assignment. Equal $(I,j)$
pairs retrieve equal inputs and replies inductively. New pairs therefore
receive independent table blocks, even when $I$ was chosen from prior
outputs and public-oracle computation.

The confirmation information lengths are $62,58,60$ bytes and each traffic
information length is $155$, so their table entries differ even if the
zero and session salts share a PRK. The marked $I^*$ is also distinct from
every earlier traffic entry. Importantly, this exact rewrite removes
the internal parent-row inputs containing a preceding secret block.
Because that row has no other interface, they leave no extra use of the
selected traffic seed. Public quantum-oracle computation is unchanged.
\end{proof}

An available test uses the public view, event ledger, declared locally
computable outputs and its independent test randomness. It cannot inspect
the designated private roots, PRKs, seeds, message key or tables. For the
application experiment below we can fix either challenge bit and then
average; the test is not given a hidden primitive key.

Let $G_R$ be this real schedule with uniform root. In $G_{\rm leaf}$ replace
\eqref{eq:joint-root-uses} by $V(\Gamma(\cdot))$ and $U(\cdot)$, replace each
column-PRK Expand interface by Lemma~\ref{lem:joint-expand-blocks}'s block
table, and replace \eqref{eq:joint-chain-step} by two independent cached
blocks $s',K$. Real confirmation/ACK operations and successor-root
derivations continue from the known replies; all signatures, checks and
later outputs are regenerated within the current game.

\begin{theorem}[A selected message leaf with joint continuation]
\label{thm:joint-message-leaf}
Let $q_R,q_0,q_1,q_s,q_M$ bound the entire ordinary compression-query work
of the five reductions below. Set
\begin{equation}
\begin{split}
 \Delta_{\rm sched}=\min\{1,\;&
 \varepsilon_{\rm Hjoint}(2,t_+,q_R,32)\\
 &+\varepsilon_{\rm Hrow}(3+4M,q_0,188)
  +\varepsilon_{\rm Hrow}(4M,q_1,188)\\
 &+\varepsilon_{\rm Hcol}(q_s)
  +\varepsilon_{\rm Hrow}(2,q_M,89)\}.
\end{split}
 \label{eq:joint-schedule-error}
\end{equation}
For every available terminal bit, the probability gap between $G_R$ and
$G_{\rm leaf}$ is at most $\Delta_{\rm sched}$.
For a common available score in $[-1,1]$, the expectation gap is at most
$2\Delta_{\rm sched}$. The second PRK-row term is omitted when the interface
itself rules out a second class.
\end{theorem}
\begin{proof}
First give every current-root column and row call to the joint primitive
experiment. Their replies suffice to compute all Expands, confirmations,
signatures, traffic and successor roots. In particular, disclosed successor
contributions and roots require no raw current root. This gives the first
term and retains the same public oracle throughout.

Next replace the at most two column-PRK rows in class-allocation order.
Preallocate the zero-salt class privately if necessary; a second class
exists only when $\Gamma(s)\ne\Gamma(0)$. For a designated class, install
the row primitive's hidden uniform key as its first $V$ reply, withheld
until allocation. The identifier is computed from known salts and $h$;
it does not inspect the private reply. Every equal $\Gamma$ reuses this
handle. Answer Expands by querying the challenge row on the actual
successive messages. Their output blocks are known, so real MACs,
verifications, signatures and continuation can all be performed.
Other class keys are sampled locally. All their real HMAC computations,
and all preparation and future work, count as ordinary queries.

Three one-block confirmation Expands and at most $M$ four-block traffic
Expands give input lengths $63,59,61$ and $156,188,188,188$. Thus the
first class has at most $3+4M$ distinct row calls and a second class at most
$4M$. A never-allocated class causes no challenge call. Their total is at
most $3+4M$; the separate worst-case budgets avoid conditioning a
primitive bound on the oracle-dependent class-equality event.

Independent $V$ replies may coincide numerically. No inequality is assumed:
known-key computations are performed through the same $h$, including
accidental overlap with the challenge key, and are counted by the
primitive experiment. An ideal row is the declared hybrid replacement.
Intentional $\Gamma$ aliases share that row in every game.

Apply Lemma~\ref{lem:joint-expand-blocks} exactly. The selected traffic seed
is now a fresh independent block, and its parent no longer exposes an
input containing it. Presample the single-column challenge seed and
install it at the newly marked entry before its blocks are produced.
There is no earlier equal-info allocation; later aliases retain this
protected handle.
Its only raw use is $F(0,s^*)$; every other block is sampled locally.
The next hop therefore makes $P_M$ independently uniform at cost
$\varepsilon_{\rm Hcol}(q_s)$.

Finally install the row challenge key at that fresh $P_M$. The two actual
Expand inputs have lengths $57$ and $89$, giving the final term.
The same exact block-table rewrite produces independent $s',K$.
The next seed may be disclosed and used in all later chain computation;
the reduction has every value it needs except the key reserved to the
application interface. These five comparisons prove the bit bound.
Randomizing a terminal bit with probability $(Z+1)/2$ proves the score
bound. Equivalently one may fix an independent application challenge bit,
apply this bound, and average over that bit.
\end{proof}

The conclusion concerns independence \emph{at key installation} with the
continuation regenerated consistently. It does not assert independence
of the key and a subsequently computed ciphertext, nor pseudorandomness
of every intermediate root relative to a fixed complete future view.

\paragraph{Application confidentiality.}
Use a one-key, one-encryption left-or-right AEAD experiment with nonce
$0^{96}$, equal-length messages, common associated data and the same
additional public oracle $h$. Its classical decryption interface returns
real plaintext or failure, except that an exact challenge query returns
an authenticated marker. Let $\varepsilon_{\rm AEAD}$ be its absolute
guessing advantage over $1/2$. An encryption-only profile uses the
corresponding one-use privacy experiment.

The honest application creates the selected send. The network adversary
receives ciphertext, metadata, status, ACKs and real authentication/error
outcomes, and can request plaintext-free $\operatorname{Resume}(id)$.
The actual payload and its retained send-intent digest are opaque
application handles. Delivery and any payload-bearing resolution output
are private to the honest application; the public continuation has no
additional payload-dependent return. An exact challenge reception can
therefore update public state using its authenticated marker and known
metadata. Other ciphertexts use the declared decryption interface.

The send-origin normalization caches the encryption reply at its first
\emph{private computation}, even if that computation precedes durable
commit or network release. Native crash recovery may repeat encryption,
but always on the same original tuple; the reduction retrieves its cache.
The cache includes the marked material's original allocation and
authoritative owner/history, direction and index. Distinct material
allocations are not merged merely because their epoch numbers agree.

\begin{corollary}[First message of a recovered direction]
\label{cor:joint-first-message}
For an event $A$ computable from the available view and ledger without
reading the hidden AEAD challenge bit $b$, let
$Z=\mathbf1_A(2\mathbf1[g=b]-1)$.
Under the application interface above,
\[
 |\mathbb E_{G_R}Z|
 \le 2\Delta_{\rm sched}+2\varepsilon_{\rm AEAD}.
\]
\end{corollary}
\begin{proof}
Theorem~\ref{thm:joint-message-leaf}, applied for each fixed $b$ and
averaged, changes this score by at most $2\Delta_{\rm sched}$.
In $G_{\rm leaf}$ give the selected encryption/decryption operations
to the AEAD challenger and sample every other block locally.
Confirmation, ACK, successor-root and subsequent-chain computation uses
known material and the same public oracle. Cache the first encryption
response; exact challenge reception needs only the marker and metadata.
At the end return $g$ when $A$ holds and a fresh fair guess otherwise.
Its centered success is exactly $\mathbb E Z$, bounded by
$2\varepsilon_{\rm AEAD}$.
\end{proof}

The public application boundary is substantive. For two challenge
plaintexts with different native send-intent digests, publishing the
result of $\operatorname{send\_message}(id,m_0,ad)$ after the challenge
would distinguish a retained result from an intent conflict. No
record acquisition is needed for that test. The network profile uses
Resume while the honest application retains the plaintext-bearing API;
original eligibility is unchanged.

\paragraph{Entry from a recovered contribution.}
The uniform value above is the epoch root \emph{after} root derivation.
For a separately justified uniform combined contribution $S$, known
previous root $r$ and public response core, the actual derivation is
\[
\begin{aligned}
 P&=F(r,S),\qquad R=F(P,I_R\Vert\mathtt{01}),\\
 I_R&=D\Vert\texttt{pending-root/HKDF-SHA256/}\\
    &\quad\Vert\operatorname{hash}(\texttt{response-core},\mathit{core}).
\end{aligned}
\]
Its information string has 97 bytes and its sole Expand message has 98.
Suppose this is one original Extract/Expand, with exact repeats cached,
and $S,P$ have no additional raw-use or disclosure channel.
The retained root-initialization lemma applies with the entire controller
above as continuation, at bit cost $\Delta_{\rm init}$. Its ideal reply
can be presampled and withheld, providing precisely the joint uniform
root initialization required here. It regenerates confirmations,
signatures, traffic and future epochs; it does not replace a root inside
a fixed public transcript.

Let $\mathsf W_{\rm cont}^{[j]}$ denote this uniform-combined-contribution
continuation, with the same candidate-local score code $\Psi_j$, and put
$b_j^{\rm cont}=\mathbb E_{\mathsf W_{\rm cont}^{[j]}}\Psi_j$.
Consequently,
\[
 |b_j^{\rm cont}|\le
 2\Delta_{{\rm init},j}
 +2\Delta_{{\rm sched},j}
 +2\varepsilon_{{\rm AEAD},j}.
\]
The native accounting term
$b_j=\mathbb E_{\widetilde{\mathsf V}^{[j]}}\Psi_j$ belongs instead to
the random-component game, which may still use the real ContextBound
combiner. If its actual complete comparison to
$\mathsf W_{\rm cont}^{[j]}$ has bit gap at most
$\varepsilon_{{\rm CB},j}$, then
\[
 |b_j-b_j^{\rm cont}|\le2\varepsilon_{{\rm CB},j},
 \qquad
 |b_j|\le2\varepsilon_{{\rm CB},j}+|b_j^{\rm cont}|.
\]
This comparison retains the controller, query interfaces and score/event
code, counts all shared-oracle preparation and continuation, and recomputes
events within each game's execution. Its actual tuple and disclosure
conditions must hold. The randomized test with success probability
$(\Psi_j+1)/2$ gives the first inequality; the second is the triangle
inequality. Native correspondence $\rho_j$, component and authentication
terms, and original-game projection residuals remain in their established
locations. This supplies the continuation bound for the specified profile,
including known old roots and allowed later disclosure.
